\chapter{Three Primes}
\begin{center}\itshape
On 7 June 1742 Christian Goldbach wrote to Euler that every number\\
seems to be a sum of three primes. Two centuries later Vinogradov\\
proved it for every large odd number.
\end{center}

Goldbach's letter of 1742 asked, in the language of the time, whether every integer greater than $2$ is a sum of three primes (Goldbach counted $1$ as a prime). Euler replied that it would follow from the statement that every even number is a sum of two primes, which he was sure of and could not prove. Nobody can, still.

Volume~5 counts partitions and sums of powers by the circle method. In 1923 Hardy and Littlewood treated three primes the same way, but only under the generalized Riemann hypothesis. In 1937 Vinogradov removed the hypothesis with his estimate of Chapter~5.

This chapter proves his theorem. The major arcs come from Siegel--Walfisz (Lemma~\ref{lem:majorarc}) and the minor arcs from Vinogradov (Theorem~\ref{thm:vinogradov3}).

\begin{reminder}[from Chapter~5]
\emph{Geometric sums.} For an interval $I$, $\bigl|\sum_{n\in I}e(\beta n)\bigr|\le\min\bigl(|I|+1,\frac1{2\|\beta\|}\bigr)$.
\emph{Primes} (Chapter~5). Dirichlet's approximation (Lemma~\ref{lem:dirichlet}); Vinogradov's estimate (Theorem~\ref{thm:vinogradov3}): if $\bigl|\alpha-\frac aq\bigr|\le\frac1{q^2}$ with $(a,q)=1$, then $S(\alpha)=\sum_{n\le N}\Lambda(n)e(n\alpha)\ll\bigl(\frac N{\sqrt q}+N^{4/5}+\sqrt{Nq}\bigr)(\log N)^5$; and the major arc lemma (Lemma~\ref{lem:majorarc}), with Ramanujan's sum $c_q(a)=\mu(q)$ for $(a,q)=1$.
\end{reminder}

\section{The singular series of three primes}

The number we want is
\[
R(N)=\sum_{n_1+n_2+n_3=N}\Lambda(n_1)\Lambda(n_2)\Lambda(n_3)=\int_0^1S(\alpha)^3e(-N\alpha)\dd\alpha,
\]
where $S(\alpha)=\sum_{n\le N}\Lambda(n)e(n\alpha)$; the integral picks out the coefficient of $e(N\alpha)$. The main term will contain a product over primes that records the local obstructions. Its ingredients are Ramanujan's sums for every $a$, not only for $a$ prime to $q$.

\begin{lemma}\label{lem:ramsum}
For $q\ge1$ and an integer $N$, let $c_q(N)=\sum_{r\le q,\,(r,q)=1}e(\frac{rN}q)$. Then $c_q(N)=\sum_{d\mid(q,N)}\mu(\frac qd)d$. In particular $c_q(N)$ is a real integer, multiplicative in $q$, and $c_p(N)=p-1$ if $p\mid N$, $c_p(N)=-1$ if $p\nmid N$.
\end{lemma}
\begin{proof}
By $\sum_{d\mid(r,q)}\mu(d)=\mathbf 1_{(r,q)=1}$,
\[
c_q(N)=\sum_{d\mid q}\mu(d)\sum_{s\le q/d}e\Bigl(\frac{sN}{q/d}\Bigr)=\sum_{d\mid q}\mu(d)\,\frac qd\,\mathbf 1_{\frac qd\mid N},
\]
and with $e=\frac qd$ this is $\sum_{e\mid q,\,e\mid N}\mu(\frac qe)e$. As a function of $q$ it is the Dirichlet convolution of $\mu$ with the multiplicative function $e\mapsto e\mathbf 1_{e\mid N}$, hence multiplicative. For $q=p$ the divisors are $e=1$, giving $-1$, and $e=p$ if $p\mid N$, giving $p$.
\end{proof}

\begin{lemma}\label{lem:phi34}
For every $q\ge1$, $\varphi(q)\ge\frac1{64}q^{3/4}$. Consequently $\sum_{q>Q}\frac1{\varphi(q)^2}\le\frac{2^{14}}{\sqrt Q}$ for $Q\ge1$.
\end{lemma}
\begin{proof}
$\frac{\varphi(q)}{q^{3/4}}=\prod_{p^j\,\|\,q}p^{j/4}(1-\frac1p)\ge\prod_{p\mid q}p^{1/4}(1-\frac1p)$. For $p\ge16$ the factor is at least $2\cdot\frac12=1$, and for the six primes $p<16$ it is at least $2^{1/4}\cdot\frac12>\frac12$. So the product is at least $2^{-6}$. Then $\sum_{q>Q}\varphi(q)^{-2}\le2^{12}\sum_{q>Q}q^{-3/2}\le2^{12}(Q^{-3/2}+\int_Q^\infty t^{-3/2}\dd t)\le2^{12}\cdot\frac3{\sqrt Q}$.
\end{proof}

\begin{proposition}\label{prop:sing3}
The \emph{singular series} $\mathfrak S(N)=\sum_{q\ge1}\frac{\mu(q)c_q(N)}{\varphi(q)^3}$ converges absolutely, and
\[
\mathfrak S(N)=\prod_{p\mid N}\Bigl(1-\frac1{(p-1)^2}\Bigr)\prod_{p\nmid N}\Bigl(1+\frac1{(p-1)^3}\Bigr).
\]
If $N$ is even, $\mathfrak S(N)=0$. If $N$ is odd, $1\le\mathfrak S(N)\le4$.
\end{proposition}
\begin{proof}
Since $|c_q(N)|\le\varphi(q)$, the terms are at most $\varphi(q)^{-2}$ in absolute value, and Lemma~\ref{lem:phi34} gives absolute convergence. The term is a multiplicative function of $q$ (Lemma~\ref{lem:ramsum}), supported on squarefree $q$, so the series equals its Euler product $\prod_p(1-\frac{c_p(N)}{(p-1)^3})$, which is the product displayed. At $p=2$ the factor is $0$ if $2\mid N$ and $2$ otherwise. For odd $N$, the factors with $p\mid N$, $p\ge3$, are at least $\prod_{m\ge2}(1-\frac1{m^2})=\frac12$ in total, since the numbers $p-1$ are distinct integers $\ge2$ and the product telescopes: $\prod_{m=2}^M\frac{(m-1)(m+1)}{m^2}=\frac{M+1}{2M}$. The factors with $p\nmid N$ are at least $1$. So $\mathfrak S(N)\ge2\cdot\frac12=1$. Similarly $\mathfrak S(N)\le2\prod_{m\ge2}(1+\frac1{m^3})\le2\exp(\zeta(3)-1)<4$.
\end{proof}

\section{Vinogradov's three primes theorem}

\begin{theorem}[Vinogradov, 1937]\label{thm:threeprimes}
For every $A>0$, as $N\to\infty$,
\[
R(N)=\frac12\mathfrak S(N)N^2+O_A\Bigl(\frac{N^2}{(\log N)^A}\Bigr).
\]
The implied constant cannot be computed by the proof.
\end{theorem}

\begin{proof}
Let $\mathcal L=\log N$, let $B=2A+12$ and $P=\mathcal L^B$, and assume $N$ is large.

\emph{The arcs.} For $1\le q\le P$ and $(a,q)=1$, let $\mathfrak M(q,a)$ be the set of $\alpha$ modulo $1$ with $|\alpha-\frac aq|\le\frac PN$. Two distinct fractions with denominators at most $P$ differ by at least $\frac1{P^2}>\frac{2P}N$, so these arcs are disjoint. Let $\mathfrak M$ be their union, the \emph{major arcs}, and $\mathfrak m$ the complement, the \emph{minor arcs}. The measure of $\mathfrak M$ is at most $\sum_{q\le P}q\cdot\frac{2P}N\le\frac{2P^3}N$.

\emph{Minor arcs.} Let $\alpha\in\mathfrak m$. By Lemma~\ref{lem:dirichlet} with $Q=\frac NP$ there are coprime $a$, $q$ with $q\le\frac NP$ and $|\alpha-\frac aq|\le\frac P{qN}\le\frac1{q^2}$. If $q\le P$, then $|\alpha-\frac aq|\le\frac PN$ and $\alpha\in\mathfrak M$; so $q>P$. Theorem~\ref{thm:vinogradov3} gives
\[
|S(\alpha)|\ll\Bigl(\frac N{\sqrt P}+N^{4/5}+\frac N{\sqrt P}\Bigr)\mathcal L^5\ll N\mathcal L^{5-B/2}.
\]
By Parseval, $\int_0^1|S|^2=\sum_{n\le N}\Lambda(n)^2\le\mathcal L\psi(N)\le2N\mathcal L$ (Chebyshev). So
\[
\Bigl|\int_{\mathfrak m}S(\alpha)^3e(-N\alpha)\dd\alpha\Bigr|\le\sup_{\mathfrak m}|S|\int_0^1|S|^2\ll N^2\mathcal L^{6-B/2}=N^2\mathcal L^{-A}.
\]

\emph{Major arcs.} Let $\alpha=\frac aq+\beta\in\mathfrak M(q,a)$ and $T(\beta)=\sum_{n\le N}e(\beta n)$. Lemma~\ref{lem:majorarc}, with $A'=4B+A$ in place of $A$, gives $S(\alpha)=\frac{\mu(q)}{\varphi(q)}T(\beta)+O(E)$ with $E=PN\mathcal L^{-A'}$, since $|\beta|N\le P$. Both $S(\alpha)$ and $\frac{\mu(q)}{\varphi(q)}T(\beta)$ have modulus at most $2N$, and $|x^3-y^3|\le3\max(|x|,|y|)^2|x-y|$, so on $\mathfrak M$ the cube $S^3$ differs from $(\frac{\mu(q)}{\varphi(q)}T)^3$ by $O(N^2E)$. Integrated over $\mathfrak M$ this costs $O(N^2E\cdot\frac{P^3}N)=O(N^2\mathcal L^{4B-A'})=O(N^2\mathcal L^{-A})$. What remains is, since $\mu(q)^3=\mu(q)$ and $\sum_{(a,q)=1}e(-\frac{aN}q)=c_q(-N)=c_q(N)$,
\[
\sum_{q\le P}\frac{\mu(q)c_q(N)}{\varphi(q)^3}\int_{-P/N}^{P/N}T(\beta)^3e(-N\beta)\dd\beta.
\]

\emph{The singular integral.} Over the full interval, $\int_{-1/2}^{1/2}T(\beta)^3e(-N\beta)\dd\beta$ counts the triples of integers in $[1,N]$ with sum $N$, which is $\binom{N-1}2=\frac{N^2}2+O(N)$. For $\frac PN\le|\beta|\le\frac12$, $|T(\beta)|\le\frac1{2|\beta|}$, so the two tails together are at most $2\int_{P/N}^\infty\frac{\dd\beta}{8\beta^3}=\frac{N^2}{8P^2}$. So the integral over $|\beta|\le\frac PN$ is $\frac{N^2}2+O(\frac{N^2}{P^2})$.

\emph{The singular series.} The sum over $q\le P$ of $|\frac{\mu(q)c_q(N)}{\varphi(q)^3}|\le\varphi(q)^{-2}$ is bounded, and by Lemma~\ref{lem:phi34} it differs from $\mathfrak S(N)$ by $O(P^{-1/2})$. So the major arcs give $\frac12\mathfrak S(N)N^2+O(N^2P^{-1/2})+O(N^2\mathcal L^{-A})$, and $P^{-1/2}=\mathcal L^{-A-6}$.
\end{proof}

\begin{corollary}[Vinogradov, 1937]\label{cor:threeprimes}
Every sufficiently large odd number is a sum of three primes.
\end{corollary}
\begin{proof}
Let $N$ be odd and large. By Proposition~\ref{prop:sing3} and the theorem, $R(N)\ge\frac12N^2-O(\frac{N^2}{\log N})\ge\frac13N^2$. The triples in which some $n_i$ is a power $p^j$ with $j\ge2$ contribute little: there are at most $\sqrt N\log_2N$ such $n_i\le N$, at most $N$ choices for the other two entries once one entry is fixed, three positions, and weight at most $(\log N)^3$, in total $O(N^{3/2}(\log N)^4)$. So the triples of primes contribute at least $\frac14N^2>0$, and some $p_1+p_2+p_3=N$.
\end{proof}

The corollary is ineffective: ``sufficiently large'' depends on the Siegel constant, which no one can compute. Later versions of the proof, which avoid Siegel's theorem by treating a possible exceptional zero separately, gave explicit bounds; Borozdin's bound of 1956, $N\ge3^{3^{15}}$, is a number with almost seven million digits. In 2013 Harald Helfgott closed the gap: every odd number greater than $5$ is a sum of three primes. His proof sharpens both arcs until the analytic argument covers every $N$ beyond a modest explicit bound, and a computer check by Helfgott and Platt, built on the verification of the binary conjecture up to $4\cdot10^{18}$ by Oliveira e Silva, Herzog and Pardi, does the rest.

\begin{figure}[ht]
\centering
\fig{ch07-threeprimes}
\caption{$R(N)$ divided by $\frac12\mathfrak S(N)N^2$, for the odd $N=1001,2001,\dots,99001$, with $R(N)$ computed exactly by a fast Fourier transform of $\Lambda$. The ratio is $0.945$ at $N=1001$, $0.993$ at $N=10001$ and $0.998$ at $N=99001$. The singular series itself moves a lot, between $1.33$ and $2.3$ for these $N$ (it is smallest when $N$ has many small prime factors), and the main term follows it.}
\label{fig:threeprimes}
\end{figure}

\begin{think}
Why does the same method fail for two primes? The minor arc bound used $\int|S|^2\approx N\log N$ and a saving of a power of $\log N$ in $\sup_{\mathfrak m}|S|$. For two primes, the minor arc integral $\int_{\mathfrak m}S(\alpha)^2e(-N\alpha)\dd\alpha$ must be compared with a main term of size $N$, and the only general bound is $\int_{\mathfrak m}|S|^2$, which is about $N\log N$ (Exercise~\ref{ex:binary}). One factor of $S$ less, and the method has nothing left to absorb the loss. What it can still prove is that \emph{almost all} even numbers are sums of two primes (Exercise~\ref{ex:almostall}).
\end{think}

\begin{beauty}
The count $R(N)$ is an integral over the circle. Near the fractions with small denominators the integrand is a product of a local factor, Ramanujan's sum modulo $q$, and an archimedean integral. Summing over $q$ assembles the local factors into the singular series, a product over the primes. The minor arcs are smaller than the main term, and that comparison is the work of the chapter.
\end{beauty}

\section{Exercises}

\begin{exercise}[$\star$]\label{ex:sing3}
Show that for odd $N$, $\mathfrak S(N)=2\prod_{p\mid N,\,p>2}\bigl(1-\frac1{(p-1)^2}\bigr)\prod_{p\nmid N,\,p>2}\bigl(1+\frac1{(p-1)^3}\bigr)$. By what factor does the main term of $R(N)$ drop when $3\mid N$, compared with a nearby $N$ prime to $3$ and with the same other prime factors? Check it against the values of $\mathfrak S(N)$ in Figure~\ref{fig:threeprimes}.
\end{exercise}
\begin{exercise}[$\star\star$]\label{ex:r3count}
Let $r_3(N)$ be the number of ordered triples of primes with $p_1+p_2+p_3=N$. Deduce from Theorem~\ref{thm:threeprimes} that $r_3(N)\sim\frac{\mathfrak S(N)N^2}{2(\log N)^3}$ for odd $N\to\infty$.
\end{exercise}
\begin{exercise}[$\star\star$]\label{ex:binary}
With the arcs of the proof of Theorem~\ref{thm:threeprimes}, show that $\int_{\mathfrak M}|S|^2\ll N\log\log N$, and deduce that $\int_{\mathfrak m}|S|^2\sim N\log N$.
\end{exercise}
\begin{exercise}[$\star\star\star$]\label{ex:almostall}
\emph{Almost all even numbers are sums of two primes} (Chudakov, van der Corput, Estermann, 1937--1938). Let $R_2(n)=\sum_{m_1+m_2=n}\Lambda(m_1)\Lambda(m_2)$. Show that for every $A>0$, all but $O_A(N(\log N)^{-A})$ even $n\in\left]\frac N2,N\right]$ have $R_2(n)\ge\frac N8$, and deduce that the number of even $n\le N$ that are not sums of two primes is $O_A(N(\log N)^{-A})$.
\end{exercise}
\section{Hints}

\begin{hint}{ex:sing3}
Read off the factor at $p=2$ and at $p=3$ in Proposition~\ref{prop:sing3}.
\end{hint}
\begin{hint}{ex:r3count}
Each triple carries a weight $w$ with $w\le(\log N)^3$, and $w\ge(1-\eta)^3(\log N)^3$ unless some $p_i\le N^{1-\eta}$.
\end{hint}
\begin{hint}{ex:binary}
On $\mathfrak M(q,a)$, $|S|\le\frac{|T|}{\varphi(q)}+O(E)$, and $\int_{-1/2}^{1/2}|T|^2\le N$.
\end{hint}
\begin{hint}{ex:almostall}
Split $R_2(n)=\int_{\mathfrak M}+\int_{\mathfrak m}$. The major arcs give $\mathfrak S_2(n;P)\,n+O(N\log P/P)+O(N\mathcal L^{-A})$ with $\mathfrak S_2(n;P)=\sum_{q\le P}\frac{\mu(q)^2c_q(n)}{\varphi(q)^2}$, which is at least $1$ unless $\tau(n)$ is unusually large. The minor arc parts $E(n)$ satisfy $\sum_n|E(n)|^2\le\int_{\mathfrak m}|S|^4$ by Parseval.
\end{hint}
\section{Solutions}

\begin{solution}{ex:sing3}
\idea Each prime contributes its own factor, so changing the divisibility by $3$ changes one factor only.
\writeup For odd $N$ the factor at $p=2$ is $1+\frac1{(2-1)^3}=2$, and the other factors are as in Proposition~\ref{prop:sing3}. The factor at $3$ is $1-\frac14=\frac34$ if $3\mid N$ and $1+\frac18=\frac98$ if not; their ratio is $\frac{3/4}{9/8}=\frac23$. So, other things equal, odd multiples of $3$ have two thirds as many representations as three primes: if $3\mid N$, a representation with all $p_i\ne3$ needs residues modulo $3$ adding to $0$, from $\{1,2\}^3$ only $(1,1,1)$ and $(2,2,2)$ work, $2$ out of $8$; if $N\equiv1$ or $2$, three out of eight work. In Figure~\ref{fig:threeprimes}, the prime $N=3001$ has $\mathfrak S(N)=2.301$, the generic value $2\prod_{p>2}(1+\frac1{(p-1)^3})$ up to a factor $1-O(\frac1{N^3})$, while $N=2001=3\cdot23\cdot29$ has $\mathfrak S(N)=1.529$; the factors at $23$ and $29$ change it by less than $0.4\%$, and $\frac{1.529}{2.301}=0.664\approx\frac23$.
\end{solution}
\begin{solution}{ex:r3count}
\idea The logarithmic weights are essentially constant, $\log N$, for all but a negligible set of triples.
\writeup Let
\[
R'(N)=\sum_{p_1+p_2+p_3=N}\log p_1\log p_2\log p_3.
\]
The proof of Corollary~\ref{cor:threeprimes} shows $R'(N)=R(N)+O(N^{3/2}(\log N)^4)=\frac12\mathfrak S(N)N^2(1+o(1))$, using $\mathfrak S(N)\ge1$. First, $R'(N)\le(\log N)^3r_3(N)$, so $r_3(N)\ge(1+o(1))\frac{\mathfrak S(N)N^2}{2(\log N)^3}$. Next fix $\eta\in\left]0,1\right[$. The triples with some $p_i\le N^{1-\eta}$ number at most $3N^{1-\eta}\cdot N$, since $p_i$ and one other prime determine the third. For the others the weight is at least $(1-\eta)^3(\log N)^3$, so
\[
r_3(N)\le\frac{R'(N)}{(1-\eta)^3(\log N)^3}+3N^{2-\eta}=\frac{(1+o(1))\mathfrak S(N)N^2}{2(1-\eta)^3(\log N)^3}.
\]
As $\eta$ is arbitrary and $\mathfrak S(N)\le4$, the asymptotic follows.
\end{solution}
\begin{solution}{ex:binary}
\idea On the major arcs $S$ is a scaled copy of $T$, whose square integral is only $N$; the factor $\frac1{\varphi(q)^2}$ summed against the $\varphi(q)$ arcs gives $\sum\frac1{\varphi(q)}\asymp\log P$.
\writeup On $\mathfrak M(q,a)$, $|S|^2\le2\frac{|T|^2}{\varphi(q)^2}+2|S-\frac{\mu(q)}{\varphi(q)}T|^2\le2\frac{|T(\beta)|^2}{\varphi(q)^2}+O(E^2)$ with $E=PN\mathcal L^{-A'}$. Since $\int_{-1/2}^{1/2}|T|^2=\lfloor N\rfloor$,
\[
\int_{\mathfrak M}|S|^2\le2N\sum_{q\le P}\frac{\varphi(q)}{\varphi(q)^2}+O\Bigl(E^2\frac{P^3}N\Bigr)=2N\sum_{q\le P}\frac1{\varphi(q)}+O(NP^5\mathcal L^{-2A'}).
\]
By Lemma~\ref{lem:phisum}, $\sum_{q\le P}\frac1{\varphi(q)}\le e(1+\log P)\ll\log\mathcal L$. With $A'$ large the error is $O(N)$, so $\int_{\mathfrak M}|S|^2\ll N\log\log N$. Since $\int_0^1|S|^2=\sum_{n\le N}\Lambda(n)^2=N\log N(1+o(1))$ (partial summation from $\psi(N)\sim N$), $\int_{\mathfrak m}|S|^2=N\log N(1+o(1))$. For two primes the main term is of order $N$, smaller than this by a factor $\log N$, so bounding $|\int_{\mathfrak m}S^2e(-N\alpha)|$ by $\int_{\mathfrak m}|S|^2$ cannot work.
\end{solution}
\begin{solution}{ex:almostall}
\idea Bessel's inequality turns a bound for $\int_{\mathfrak m}|S|^4$, which is fine (one factor of $\sup|S|$ to spare), into a bound for the number of $n$ whose minor arc contribution is large.
\writeup Use the arcs of Theorem~\ref{thm:threeprimes} with $P=\mathcal L^B$, $B$ to be chosen, and $A'$ large. For $n\le N$, $R_2(n)=\int_0^1S^2e(-n\alpha)=M(n)+E(n)$, the major and minor parts.
\emph{Minor arcs.} $E(n)$ is the $n$-th Fourier coefficient of $S^2\mathbf 1_{\mathfrak m}$, so by Parseval $\sum_{n}|E(n)|^2=\int_{\mathfrak m}|S|^4\le\sup_{\mathfrak m}|S|^2\int_0^1|S|^2\ll N^2\mathcal L^{10-B}\cdot N\mathcal L$. Hence $|E(n)|\ge\frac N8$ for at most $O(N\mathcal L^{11-B})$ values of $n$.
\emph{Major arcs.} As in the theorem, $|S^2-\frac{\mu(q)^2}{\varphi(q)^2}T^2|\ll NE$ on $\mathfrak M(q,a)$, which costs $O(NE\frac{P^3}N)=O(N\mathcal L^{4B-A'})$, and
\[
M(n)=\sum_{q\le P}\frac{\mu(q)^2c_q(n)}{\varphi(q)^2}\int_{-P/N}^{P/N}T^2e(-n\beta)\dd\beta+O(N\mathcal L^{4B-A'}).
\]
The full integral is $n-1$ (pairs in $[1,N]^2$ with sum $n\le N$), and the tails are at most $2\int_{P/N}^\infty\frac{\dd\beta}{4\beta^2}=\frac N{2P}$. As $|c_q(n)|\le\varphi(q)$, $\sum_{q\le P}\frac{|c_q(n)|}{\varphi(q)^2}\le e(1+\log P)$, so $M(n)=\mathfrak S_2(n;P)(n-1)+O(\frac{N\log P}P)+O(N\mathcal L^{4B-A'})$.
\emph{The singular series.} As in Proposition~\ref{prop:sing3}, the full series $\mathfrak S_2(n)=\sum_q\frac{\mu(q)^2c_q(n)}{\varphi(q)^2}$, if absolutely convergent, equals $\prod_{p\mid n}(1+\frac1{p-1})\prod_{p\nmid n}(1-\frac1{(p-1)^2})$, which for even $n$ is $2\prod_{p\mid n,p>2}\frac p{p-1}\prod_{p\nmid n,p>2}(1-\frac1{(p-1)^2})\ge2\cdot\frac12=1$, by the telescoping product. For convergence and the tail, by Lemma~\ref{lem:ramsum} $|c_q(n)|\le\sum_{d\mid(q,n)}d$, so, with Lemma~\ref{lem:phi34},
\[
\sum_{q>P}\frac{\mu(q)^2|c_q(n)|}{\varphi(q)^2}\le2^{12}\sum_{d\mid n}d\sum_{m>P/d}(dm)^{-3/2}\le2^{12}\sum_{d\mid n}d\cdot d^{-3/2}\cdot3\Bigl(\frac dP\Bigr)^{1/2}=\frac{3\cdot2^{12}\tau(n)}{\sqrt P},
\]
using $\sum_{m>Y}m^{-3/2}\le3Y^{-1/2}$ for $Y\ge1$ and the trivial bound $\sum_mm^{-3/2}\le3\le3(\frac dP)^{1/2}$ when $d>P$. The $n\le N$ with $\tau(n)>P^{1/4}$ number at most $P^{-1/4}\sum_{n\le N}\tau(n)\le2N\mathcal L^{1-B/4}$. For the other even $n$, $\mathfrak S_2(n;P)\ge1-O(P^{-1/4})\ge\frac34$.
\emph{Conclusion.} For even $n\in\left]\frac N2,N\right]$ outside both exceptional sets, $R_2(n)\ge\frac34(\frac N2-1)-o(N)-\frac N8\ge\frac N8$ for $N$ large. Choosing $B=4A+20$ and $A'=4B+A$, the exceptional sets have $O(N\mathcal L^{-A})$ elements. When $R_2(n)\ge\frac N8$, the pairs involving a prime power $p^j$, $j\ge2$, contribute $O(\sqrt N\mathcal L^3)$, so $n$ is a sum of two primes. Applying this to $N,\frac N2,\frac N4,\dots$ covers all even $n\le N$ with $O(N\mathcal L^{-A})$ exceptions.
\end{solution}
\section*{Sources and further reading}

\begin{itemize}
\item \emph{History.} Goldbach's letter to Euler of 7 June 1742 and Euler's reply of 30 June 1742 are in P.-H.~Fuss, \emph{Correspondance mathématique et physique de quelques célèbres géomètres du XVIII\textsuperscript{ème} siècle}, St.\ Petersburg, 1843, vol.~1.
\item \emph{The circle method.} G.~H.~Hardy and J.~E.~Littlewood, ``Some problems of `Partitio Numerorum' III: On the expression of a number as a sum of primes'', \emph{Acta Math.}\ 44 (1923), 1--70.
\item \emph{Three primes.} I.~M.~Vinogradov, ``Representation of an odd number as a sum of three primes'', \emph{C.\ R.\ (Doklady) Acad.\ Sci.\ URSS} 15 (1937), 291--294. K.~G.~Borozdin, an explicit bound announced in 1956; see the survey in H.~A.~Helfgott, ``The ternary Goldbach problem'', arXiv:1501.05438. H.~A.~Helfgott, ``The ternary Goldbach conjecture is true'', arXiv:1312.7748 (2013). The almost-all result: N.~G.~Chudakov, \emph{Dokl.\ Akad.\ Nauk SSSR} 17 (1937); J.~G.~van der Corput, \emph{Acta Arith.}\ 2 (1937); T.~Estermann, \emph{Proc.\ London Math.\ Soc.}\ (2) 44 (1938).
\end{itemize}
